\documentclass{article}

\usepackage[final]{neurips_2023}
\usepackage[utf8]{inputenc}
\usepackage[T1]{fontenc}
\usepackage{hyperref}
\usepackage{url}
\usepackage{booktabs}
\usepackage{amsfonts}
\usepackage{nicefrac}
\usepackage{microtype}
\usepackage{xcolor}
\usepackage{graphicx}
\usepackage{subcaption}
\usepackage{multirow}
\usepackage{siunitx}
\usepackage{enumitem}
\usepackage{array}
\usepackage{tabularx}
\usepackage{placeins}
\usepackage{float}

\title{AMAM-128: A Compact Multi-Regime Benchmark for Metallography Segmentation}

\author{%
  Anonymous Authors\\
  Anonymous Institution\\
  \texttt{anonymous@anonymous.edu}
}

\newcommand{\datasetname}{AMAM-128}
\newcommand{\datasetfull}{Annotated Metallic Alloys Microstructures}
\newcommand{\paircount}{128}
\newcommand{\subsetcount}{6}
\newcommand{\origcount}{136}
\newcommand{\paircoverage}{94.1\%}
\newcommand{\tablefontsize}{\small}

\begin{document}

\maketitle

\begin{abstract}
Metallography segmentation benchmarks often overestimate progress because semantics, magnification, and artifact profiles shift across material systems. Existing datasets are commonly strong in one regime but weak for cross-regime stress testing. We present \datasetfull{} (\datasetname{}), a strict pair-only benchmark release with \paircount{} original--label tuples across \subsetcount{} curated subsets. Unlike ad-hoc photo collections, \datasetname{} is constructed along explicit benchmarking axes (material family, processing condition, magnification regime, artifact profile) while preserving subset-native phase semantics, so each subset acts as a controlled stress slice rather than a random image pool. We report three benchmark layers: a 10-method classical survey, a 29-model supervised deep survey (14 general + 15 metallography-oriented, including an MLography-style U-Net variant), and a 6-model foundation/edge add-on survey (SAM, SlimSAM, TextureSAM, HED, PidiNet). The central empirical result is that all evaluated families still underperform on this benchmark (best macro mIoU 0.6519; gap-to-perfect 0.3481), confirming substantial open headroom for new methods. This multi-layer, multi-objective design therefore provides a representative stress benchmark for in-domain quality, cross-material transfer, cross-scale transfer, and artifact robustness. Benchmark website (interactive explorer and downloadable artifacts): \url{https://scientific-computing-user.github.io/amam-benchmark-site/}.
\end{abstract}

\section{Introduction}
Automated microstructure segmentation is foundational for quantitative metallography, yet benchmark methodology remains underdeveloped. The main challenge is not only annotation count: alloy family, process state, magnification, and imaging artifacts all change the segmentation regime. Benchmarks that flatten these shifts into one pooled setting can overstate progress.

Prior resources enabled important progress, including UHCS-focused datasets, broader materials-segmentation efforts, and metallography-specific pipelines~\citep{decost2017uhcsdb,decost2019high,decost2020scientific,luengo2022tutorial,rusanovsky2020mlography,rusanovsky2022end}. However, many public settings remain narrow in at least one benchmark-critical axis: material-family breadth, cross-scale coverage, artifact-stress auditing, or strict split protocol. Recent foundation segmentation systems also transfer well in natural imagery but remain fragile on texture-heavy metallography boundaries and artifact confounds~\citep{kirillov2023segment,ji2023segment,ma2024segment,rusanovsky2023universal,rusanovsky2025texturesam}. A compact landscape comparison is given in Appendix~B (Table~\ref{tab:benchmark_landscape}).

\datasetfull{} (\datasetname{}) addresses this gap with a strict rule: only image--mask pairs are included. The release contains \paircount{} paired tuples from \origcount{} local originals (\paircoverage{} pairing rate), with no unpaired samples in evaluation. It spans two material families, four process conditions, and six magnification regimes while preserving subset-native phase semantics.

Our contribution is benchmark-centric: (i) strict pair-only release specification with deterministic splits, (ii) multi-regime design coupling material/process/scale/artifact stressors, and (iii) a 45-method audit (10 classical, 29 supervised deep, 6 foundation/edge) showing persistent underperformance (best macro mIoU 0.6519). This makes the benchmarking gap measurable rather than anecdotal.

\section{The AMAM-128 Benchmark}
\textbf{Benchmark design.} \datasetname{} includes six subsets: 4130 steel (standard), 5884 armor steel (mixed magnification), 418 printed 17-4 PH (x5 and x20), and 6280 gray cast iron (low and high magnification). Label spaces remain subset-specific to avoid invalid class merging across materially different phase taxonomies. Figure~\ref{fig:workflow} summarizes the annotation/refinement pipeline used to construct the final paired release.

\begin{figure}[t]
    \centering
    \includegraphics[width=\linewidth,trim=25 20 45 60,clip]{figures/amam-workflow.png}
    \caption{AMAM annotation and refinement workflow. From left to right, metallographers provide representative micrographs, create sparse click-level supervision, run model-assisted dense segmentation, and iteratively correct predictions. Accepted masks are then frozen into the paired benchmark release, so each final tuple keeps human-verified phase boundaries while scaling annotation throughput.}
    \label{fig:workflow}
\end{figure}

\textbf{Why this benchmark is stronger than monolithic alternatives.} \datasetname{} is intentionally compact but heterogeneous: it combines ferrite/pearlite separation, dendrite/interdendritic structure, graphite-flake segmentation, and printed-alloy grain patterns in one suite. This exposes true regime shift instead of repetitive single-regime fitting.

\textbf{Evaluation tracks.} We recommend three tracks: (i) within-subset segmentation, (ii) cross-material transfer (steel$\leftrightarrow$cast iron), and (iii) cross-scale transfer (held-out magnification). Reporting should prioritize subset-macro mIoU (Dice/Pixel Accuracy secondary) to prevent large subsets from dominating scores.
\textbf{Release accounting and category audit.} Verified release statistics are reported in Table~\ref{tab:globalstats}, and expanded aligned image-level evidence appears in Tables~\ref{tab:rep_visual_diversity} and~\ref{tab:rep_visual_diversity_part2}.

\begin{figure}[t]
    \centering
    \includegraphics[width=0.24\linewidth]{figures/4130-steel-original.png}
    \includegraphics[width=0.24\linewidth]{figures/4130-steel-labeled.png}
    \includegraphics[width=0.24\linewidth]{figures/6280-cast-iron-low-original.jpg}
    \includegraphics[width=0.24\linewidth]{figures/6280-cast-iron-low-labeled.jpg}
    \\
    \includegraphics[width=0.24\linewidth]{figures/5884-armor-steel-original.jpg}
    \includegraphics[width=0.24\linewidth]{figures/5884-armor-steel-labeled.jpg}
    \includegraphics[width=0.24\linewidth]{figures/418-17-4ph-x5-original.jpg}
    \includegraphics[width=0.24\linewidth]{figures/418-17-4ph-x5-labeled.jpg}
    \caption{Representative AMAM-128 image--mask pairs from four distinct regimes. Top row: 4130 steel and low-magnification cast iron (original/mask pairs). Bottom row: mixed-magnification armor steel and printed 17-4 PH x5 (original/mask pairs). The diversity in texture scale, boundary sharpness, and phase morphology illustrates why regime-aware benchmarking is necessary.}
    \label{fig:main_diversity_pairs}
\end{figure}

\section{Model Audit Across Three Families}
\textbf{Benchmarking design.} We evaluate three families: (i) classical (10), (ii) supervised deep (29, general-purpose + metallography-oriented), and (iii) foundation/edge add-ons (6, inference-only). All use strict pair-only inclusion, deterministic splits/seeds, and subset-macro reporting (full protocol in Reproducibility Appendix).

\textbf{Main empirical evidence.} Figure~\ref{fig:gap_overview} summarizes all 45 methods. The best result is 0.6519 mIoU and the median is 0.5362, confirming the benchmark remains unsaturated even under broad architectural and preprocessing coverage.

\begin{figure}[t]
    \centering
    \includegraphics[width=\linewidth]{figures/benchmark-gap-overview.pdf}
    \caption{
    Cross-family benchmark gap over all 45 tested methods (10 classical, 14 deep-general, 15 deep-metallography, 6 foundation/edge). Left: per-group distributions with best-in-group markers. Right: global ranking and gap to perfect overlap (1.0). Full numeric rankings are reported in Tables~\ref{tab:classical-full}, \ref{tab:deep-general}, \ref{tab:deep-metal}, and~\ref{tab:foundation-edge}.
    }
    \label{fig:gap_overview}
\end{figure}

\textbf{Cross-family findings.} No family saturates this benchmark; metallurgy-aware deep variants help but still leave substantial error; prompt-free foundation transfer remains clearly below trained supervised models.

\section{Significance and Benchmarking Guidance}
\textbf{What we learn scientifically.} The strongest model reaches 0.6519 macro mIoU, leaving a 0.3481 gap to perfect overlap; the strongest supervised deep variant reaches 0.6427; the strongest foundation add-on reaches 0.4565. This confirms substantial headroom across all model families.

\textbf{Merged evidence table (former top-performer table + benchmark-answer table).} Table~\ref{tab:merged_main_evidence} consolidates the leading model per family together with the benchmark-level conclusions those results support.
\begin{table}[H]
    \centering
    \tablefontsize
    \caption{Merged evidence summary: top-per-family results and benchmark-level conclusions (integrating former separate summary tables).}
    \label{tab:merged_main_evidence}
    \setlength{\tabcolsep}{3pt}
    \begin{tabularx}{\linewidth}{p{0.24\linewidth}p{0.24\linewidth}p{0.09\linewidth}X}
        \toprule
        Evidence axis & Best method / statistic & mIoU & Benchmark implication \\
        \midrule
        Classical best & RF (pixel features) & 0.6519 & Strongest overall score still leaves large unsolved error mass. \\
        Deep-general best & U-Net (EfficientNet-B0) & 0.6126 & Generic supervised deep models improve consistency but remain below top classical score. \\
        Deep-metallography best & Metal-U-Net++ CLAHE (EfficientNet-B0) & 0.6427 & Domain-aware preprocessing helps, but does not close the gap to reliable high-fidelity segmentation. \\
        Foundation/edge best & SAM ViT-Base (auto-mask) & 0.4565 & Prompt-free transfer remains far from benchmark-ready performance in this regime. \\
        Global distribution & Median over 45 methods & 0.5362 & Most methods cluster well below saturation, confirming persistent benchmark difficulty. \\
        Regime-shift stress & Family + magnification holdouts & --- & Error increases under material/scale shift and artifact-heavy slices; robustness remains unresolved. \\
        \bottomrule
    \end{tabularx}
\end{table}

These outcomes show that AMAM-128 is not a solved dataset despite visually simple regions in some samples. Progress should be judged by subset-aware robustness and transfer behavior, not a single pooled score.

Detailed reporting requirements and validity notes are provided in the Reporting/Validity Appendix.

\FloatBarrier
\clearpage
\bibliographystyle{plainnat}
\bibliography{refs}

\clearpage
\appendix

\section{Related Work (Appendix)}
Metallography datasets such as UHCSDB and related quantitative studies demonstrated the value of supervised microstructure segmentation and deep-learning pipelines~\citep{decost2017uhcsdb,decost2019high}. Broader materials-segmentation discussions and dataset/tutorial efforts (e.g., MetalDAM-oriented analysis) emphasized taxonomy design and practical deployment constraints~\citep{luengo2022tutorial}. Related metallography works spanning MLography, end-to-end quantitative metallography, texture-boundary modeling, and TextureSAM-style adaptation also inform this benchmark context~\citep{rusanovsky2020mlography,rusanovsky2022end,rusanovsky2023universal,rusanovsky2025texturesam}. Foundational segmentation architectures (e.g., U-Net) remain common baselines~\citep{ronneberger2015u}, while promptable foundation models and transfer analyses continue to expand~\citep{kirillov2023segment,ji2023segment,ma2024segment}. From a benchmarking perspective, the key open need is protocol rigor: stable splits, transparent metadata, and comparable reporting for cross-regime claims~\citep{mueller2016machine,ramprasad2017machine,decost2020scientific}.

\section{External Benchmark Landscape (Appendix)}
Table~\ref{tab:benchmark_landscape} and Table~\ref{tab:benchmark_landscape_part2} provide a reorganized benchmark-facing landscape of representative external metallography datasets used in prior work. Counts denote annotated units as reported by the original source (full images or annotated crops/patches, depending on the dataset protocol).
\begin{table*}[ht]
    \centering
    \scriptsize
    \caption{Summary of quantitative metallography datasets (Part I).}
    \label{tab:benchmark_landscape}
    \setlength{\tabcolsep}{3pt}
    \renewcommand{\arraystretch}{1.05}
    \begin{tabularx}{\textwidth}{p{0.16\textwidth}p{0.08\textwidth}X X}
        \toprule
        Dataset Name & No. & Description & Paper / Reference \\
        \midrule
        MetaIDAM & 42 & SEM images of steels from additive manufacturing; additional unlabeled images are often reported alongside the labeled split. & Luengo et al., ``A tutorial on the segmentation of metallographic images: taxonomy, MetalDAM dataset, deep-learning ensemble, and challenges'' (Information Fusion, 2022). \\
        TBM & 320 & 128$\times$128 annotated crops for texture-boundary detection across microscopy/metallography textures. & ``Universal Semantic-less Texture Boundary Detection for Microscopy (and Metallography)'' (NeurIPS ML4PS 2023 workshop paper). \\
        unet-Metallography & 29 & Public GitHub collection used in U-Net metallography segmentation demonstrations. & GitHub: sidav/unet-metallography repository. \\
        Artificial & 336 & Real and synthetic grain images for grain/boundary segmentation benchmarking. & ``Grain and Grain Boundary Segmentation using Machine Learning with Real and Generated Datasets'' (2023). \\
        UHCS & 961 & UHCS-focused SEM micrographs with annotations for cementite and related microconstituents. & DeCost et al., UHCSDB and related UHCS quantitative studies (2017/2019). \\
        SGMD & 60 & Images across several alloys (including Ni-, Fe-, brass-, and steel-related examples) for grain-size segmentation benchmarking. & ``An Improved U-Net Image Segmentation Method and Its Application for Metallic Grain Size Statistics'' (2022). \\
        \bottomrule
    \end{tabularx}
\end{table*}

\begin{table*}[ht]
    \centering
    \scriptsize
    \caption{Summary of quantitative metallography datasets (Part II).}
    \label{tab:benchmark_landscape_part2}
    \setlength{\tabcolsep}{3pt}
    \renewcommand{\arraystretch}{1.05}
    \begin{tabularx}{\textwidth}{p{0.16\textwidth}p{0.08\textwidth}X X}
        \toprule
        Dataset Name & No. & Description & Paper / Reference \\
        \midrule
        LFTi64 & 40 & Optical Ti-6Al-4V micrographs (20 bimodal equiaxed, 20 lamellar) used for microstructure fingerprinting. & ``Digital Fingerprinting of Microstructures'' line of work (Ti-6Al-4V). \\
        Spheroidite & 24 & Weakly supervised segmentation setting centered on spheroidite particle delineation. & Source link provided in prior notes: arXiv:1805.08693. \\
        HY-80 & 110 & Segmented HY-80 steel images used for comparing ML segmentation against thresholding-style baselines. & Mendeley dataset: rk4vbck5n2 (HY-80 steel segmentation release). \\
        NBS2 & 12 & Ni-based superalloy (Super-2) segmentation benchmark split (matrix/precipitate). & NASA pretrained microscopy models benchmark data (Super2). \\
        NBS3 & 5 & Ni-based superalloy Super-4 test-only set used for transfer/generalization checks. & NASA pretrained microscopy models benchmark data (Super4). \\
        AZA & 42 & XCT Al-Zn-alloy grayscale micrographs of dendritic structures. & DOI: 10.1016/j.matchar.2020.110119. \\
        \bottomrule
    \end{tabularx}
\end{table*}

Benchmark implication: these resources are valuable and often technically strong, but most remain narrow along at least one benchmark-critical axis (cross-material breadth, multi-scale regime coverage, explicit artifact stress testing, or strict pair-only release governance). AMAM-128 is designed as a stress benchmark rather than only a dataset collection: it combines strict pair-only tuples with explicit material/process/magnification slicing and reproducible reporting protocol in one release.

Validation note: all rows above were rechecked against the cited source papers/repositories/dataset pages; counts are reported in the source unit convention (full images vs annotated crops/patches). The Spheroidite entry retains the source pointer provided in prior notes and should be re-audited against the final camera-ready bibliography before submission.

\section{Release Statistics (Appendix)}
Table~\ref{tab:globalstats} reports the verified global release accounting for strict pair-only AMAM-128.
\begin{table}[ht]
    \centering
    \tablefontsize
    \caption{Verified release statistics for \datasetname{} (pair-only policy).}
    \label{tab:globalstats}
    \begin{tabular}{lc}
        \toprule
        Statistic & Value \\
        \midrule
        Paired tuples & \paircount{} \\
        Local originals scanned & \origcount{} \\
        Pair coverage over locals & \paircoverage{} \\
        Included subsets & \subsetcount{} \\
        Material families & 2 (Steel, Cast Iron) \\
        Distinct phase labels (across subsets) & 9 \\
        \bottomrule
    \end{tabular}
\end{table}

\section{Subset-Level Coverage (Appendix)}
Table~\ref{tab:subsetcoverage} reports the strict pair coverage profile used in this release. The key benchmarking implication is that subset-level test sizes are intentionally uneven, so macro-averaged reporting is necessary to prevent large subsets from dominating conclusions.
\begin{table}[ht]
    \centering
    \tablefontsize
    \caption{Subset-level pair coverage in \datasetname{}.}
    \label{tab:subsetcoverage}
    \begin{tabular}{lccc}
        \toprule
        Subset & Pairs & Local Originals & Coverage \\
        \midrule
        4130 Steel (Standard) & 40 & 46 & 87\% \\
        6280 Gray Cast Iron (Low) & 11 & 11 & 100\% \\
        6280 Gray Cast Iron (High) & 10 & 10 & 100\% \\
        5884 Armor Steel (Mixed) & 21 & 22 & 95\% \\
        418 Printed 17-4 PH (x5) & 25 & 25 & 100\% \\
        418 Printed 17-4 PH (x20) & 21 & 22 & 95\% \\
        \midrule
        \textbf{Total} & \textbf{128} & \textbf{136} & \textbf{94.1\%} \\
        \bottomrule
    \end{tabular}
\end{table}

\section{OOD Split Cardinalities (Appendix)}
Table~\ref{tab:oodsplits} reports the canonical out-of-distribution split sizes used in this release for family and magnification holdout analysis.
\begin{table}[ht]
    \centering
    \tablefontsize
    \caption{Illustrative OOD split cardinalities from \datasetname{} (train/test pairs).}
    \label{tab:oodsplits}
    \begin{tabular}{lcc}
        \toprule
        Split setup & Train pairs & Test pairs \\
        \midrule
        Family holdout: train Steel, test Cast Iron & 107 & 21 \\
        Family holdout: train Cast Iron, test Steel & 21 & 107 \\
        Magnification holdout: test Standard & 88 & 40 \\
        Magnification holdout: test x5 & 103 & 25 \\
        Magnification holdout: test x20 & 107 & 21 \\
        Magnification holdout: test Low & 117 & 11 \\
        Magnification holdout: test High & 118 & 10 \\
        \bottomrule
    \end{tabular}
\end{table}

\section{Representative Visual Diversity (Appendix)}
Tables~\ref{tab:rep_visual_diversity} and~\ref{tab:rep_visual_diversity_part2} expand the representative visual audit to all six included subsets and tie each sample to explicit benchmark axes (material family, magnification, process condition, and artifact profile). For each sample, we show the aligned original, ground-truth mask, and best-per-family prediction selected from the main benchmarking tables: best classical (RF), best general supervised deep (U-Net EfficientNet-B0), and best metallography-oriented deep (Metal-U-Net++ CLAHE EfficientNet-B0).

\begin{table*}[ht]
    \centering
    \tablefontsize
    \setlength{\tabcolsep}{2pt}
    \renewcommand{\arraystretch}{1.08}
    \caption{Expanded representative visual diversity audit (Part I/II): three subsets with aligned masks and best-per-family predictions (best classical: RF; best deep: U-Net EfficientNet-B0; best metallography deep: Metal-U-Net++ CLAHE EfficientNet-B0).}
    \label{tab:rep_visual_diversity}
    \begin{tabular}{p{0.105\textwidth}p{0.14\textwidth}p{0.17\textwidth}ccccc}
        \toprule
        Sample & Category focus & Benchmark relevance & Original & Mask & Best Classic & Best Deep & Best Metal \\
        \midrule
        4130 Steel\\(Standard) &
        Steel family; normalized condition; standard scale &
        Fine ferrite/pearlite granularity with dense micro-islands probes precision under small-region consistency. &
        \includegraphics[width=0.085\textwidth]{figures/appendix_preds/4130-steel-viz-original.png} &
        \includegraphics[width=0.085\textwidth]{figures/appendix_preds/4130-steel-viz-mask.png} &
        \includegraphics[width=0.085\textwidth]{figures/appendix_preds/4130-steel-pred-classic-rf.png} &
        \includegraphics[width=0.085\textwidth]{figures/appendix_preds/4130-steel-pred-deep-unet-effb0.png} &
        \includegraphics[width=0.085\textwidth]{figures/appendix_preds/4130-steel-pred-metal-unetpp-clahe-effb0.png} \\
        \midrule
        6280 Gray Cast\\Iron (Low) &
        Cast-iron family; low magnification; as-cast &
        Coarse graphite flakes against ferrite matrix test class-imbalance and dendrite/flake boundary separation. &
        \includegraphics[width=0.085\textwidth]{figures/appendix_preds/6280-cast-iron-low-viz-original.png} &
        \includegraphics[width=0.085\textwidth]{figures/appendix_preds/6280-cast-iron-low-viz-mask.png} &
        \includegraphics[width=0.085\textwidth]{figures/appendix_preds/6280-cast-iron-low-pred-classic-rf.png} &
        \includegraphics[width=0.085\textwidth]{figures/appendix_preds/6280-cast-iron-low-pred-deep-unet-effb0.png} &
        \includegraphics[width=0.085\textwidth]{figures/appendix_preds/6280-cast-iron-low-pred-metal-unetpp-clahe-effb0.png} \\
        \midrule
        6280 Gray Cast\\Iron (High) &
        Cast-iron family; high magnification; as-cast &
        Finer local morphology at high scale stresses boundary fidelity and magnification transfer from low-scale training contexts. &
        \includegraphics[width=0.085\textwidth]{figures/appendix_preds/6280-cast-iron-high-viz-original.png} &
        \includegraphics[width=0.085\textwidth]{figures/appendix_preds/6280-cast-iron-high-viz-mask.png} &
        \includegraphics[width=0.085\textwidth]{figures/appendix_preds/6280-cast-iron-high-pred-classic-rf.png} &
        \includegraphics[width=0.085\textwidth]{figures/appendix_preds/6280-cast-iron-high-pred-deep-unet-effb0.png} &
        \includegraphics[width=0.085\textwidth]{figures/appendix_preds/6280-cast-iron-high-pred-metal-unetpp-clahe-effb0.png} \\
        \bottomrule
    \end{tabular}
\end{table*}

\begin{table*}[ht]
    \centering
    \tablefontsize
    \setlength{\tabcolsep}{2pt}
    \renewcommand{\arraystretch}{1.08}
    \caption{Expanded representative visual diversity audit (Part II/II): remaining three subsets showing mixed magnification and printed-alloy stress slices with aligned best-per-family predictions.}
    \label{tab:rep_visual_diversity_part2}
    \begin{tabular}{p{0.105\textwidth}p{0.14\textwidth}p{0.17\textwidth}ccccc}
        \toprule
        Sample & Category focus & Benchmark relevance & Original & Mask & Best Classic & Best Deep & Best Metal \\
        \midrule
        5884 Armor\\Steel (Mixed) &
        Steel family; Q\&T; mixed magnification &
        Mixed magnification and dendrite/interdendrite ambiguity expose scale transfer and uncertain boundary placement. &
        \includegraphics[width=0.085\textwidth]{figures/appendix_preds/5884-armor-steel-viz-original.png} &
        \includegraphics[width=0.085\textwidth]{figures/appendix_preds/5884-armor-steel-viz-mask.png} &
        \includegraphics[width=0.085\textwidth]{figures/appendix_preds/5884-armor-steel-pred-classic-rf.png} &
        \includegraphics[width=0.085\textwidth]{figures/appendix_preds/5884-armor-steel-pred-deep-unet-effb0.png} &
        \includegraphics[width=0.085\textwidth]{figures/appendix_preds/5884-armor-steel-pred-metal-unetpp-clahe-effb0.png} \\
        \midrule
        418 Printed\\17-4 PH (x5) &
        Steel family; as-sintered printed alloy; low scale &
        Pores and local inclusions create sparse artifacts that challenge false-positive control in phase maps. &
        \includegraphics[width=0.085\textwidth]{figures/appendix_preds/418-17-4ph-x5-viz-original.png} &
        \includegraphics[width=0.085\textwidth]{figures/appendix_preds/418-17-4ph-x5-viz-mask.png} &
        \includegraphics[width=0.085\textwidth]{figures/appendix_preds/418-17-4ph-x5-pred-classic-rf.png} &
        \includegraphics[width=0.085\textwidth]{figures/appendix_preds/418-17-4ph-x5-pred-deep-unet-effb0.png} &
        \includegraphics[width=0.085\textwidth]{figures/appendix_preds/418-17-4ph-x5-pred-metal-unetpp-clahe-effb0.png} \\
        \midrule
        418 Printed\\17-4 PH (x20) &
        Steel family; as-sintered printed alloy; high scale &
        High magnification reshapes pore visibility and local texture statistics, probing scale-specific robustness. &
        \includegraphics[width=0.085\textwidth]{figures/appendix_preds/418-17-4ph-x20-viz-original.png} &
        \includegraphics[width=0.085\textwidth]{figures/appendix_preds/418-17-4ph-x20-viz-mask.png} &
        \includegraphics[width=0.085\textwidth]{figures/appendix_preds/418-17-4ph-x20-pred-classic-rf.png} &
        \includegraphics[width=0.085\textwidth]{figures/appendix_preds/418-17-4ph-x20-pred-deep-unet-effb0.png} &
        \includegraphics[width=0.085\textwidth]{figures/appendix_preds/418-17-4ph-x20-pred-metal-unetpp-clahe-effb0.png} \\
        \bottomrule
    \end{tabular}
\end{table*}

\section{Model Survey Tables (Appendix)}
Tables~\ref{tab:classical-full}, \ref{tab:deep-general}, \ref{tab:deep-metal}, and~\ref{tab:foundation-edge} provide complete rankings for all model families reported in this paper.
\begin{table}[ht]
    \centering
    \tablefontsize
    \caption{Classical baseline models (10) on \datasetname{} (macro over subsets).}
    \label{tab:classical-full}
    \begin{tabular}{lccc}
        \toprule
        Method & mIoU & Dice & Pixel Acc \\
        \midrule
        RF (pixel features) & \textbf{0.6519} & \textbf{0.7580} & 0.8193 \\
        GMM-RGB & 0.6461 & 0.7486 & \textbf{0.8567} \\
        Linear SVM (pixel features) & 0.6401 & 0.7449 & 0.8363 \\
        Gabor-KMeans & 0.5955 & 0.7063 & 0.8114 \\
        KMeans-RGB & 0.5837 & 0.6829 & 0.8142 \\
        SLIC+KMeans & 0.5363 & 0.6534 & 0.7705 \\
        Felzenszwalb-GMM & 0.4916 & 0.5938 & 0.7624 \\
        LBP+KMeans & 0.4638 & 0.5927 & 0.6767 \\
        Sobel+Watershed & 0.3551 & 0.4925 & 0.5714 \\
        Canny+Watershed & 0.3548 & 0.4423 & 0.6703 \\
        \bottomrule
    \end{tabular}
\end{table}

\begin{table*}[ht]
    \centering
    \tablefontsize
    \caption{General supervised deep models (14) on \datasetname{} (macro over subsets).}
    \label{tab:deep-general}
    \begin{tabular}{llccc}
        \toprule
        Model & Category & mIoU & Dice & Pixel Acc \\
        \midrule
        U-Net (EfficientNet-B0) & Encoder-decoder & \textbf{0.6126} & \textbf{0.7220} & 0.8165 \\
        DeepLabV3+ (EfficientNet-B0) & Context-atrous & 0.6000 & 0.7111 & \textbf{0.8349} \\
        U-Net (ResNet34) & Encoder-decoder & 0.5853 & 0.6911 & 0.8184 \\
        MAnet (ResNet34) & Attention & 0.5794 & 0.6924 & 0.7606 \\
        SegFormer (MiT-B2) & Transformer & 0.5781 & 0.6850 & 0.8009 \\
        DeepLabV3+ (ResNet34) & Context-atrous & 0.5612 & 0.6818 & 0.7734 \\
        U-Net++ (ResNet34) & Encoder-decoder & 0.5516 & 0.6666 & 0.7663 \\
        SegFormer (MiT-B0) & Transformer & 0.5451 & 0.6560 & 0.7740 \\
        PSPNet (ResNet34) & Pyramid & 0.4818 & 0.5912 & 0.7352 \\
        FPN (ResNet34) & Pyramid & 0.4598 & 0.5750 & 0.6940 \\
        UPerNet (MiT-B0) & Transformer & 0.4501 & 0.5539 & 0.6486 \\
        LinkNet (ResNet34) & Lightweight & 0.4343 & 0.5550 & 0.6561 \\
        UPerNet (MiT-B2) & Transformer & 0.3878 & 0.4869 & 0.5977 \\
        DeepLabV3 (ResNet34) & Context-atrous & 0.3430 & 0.4134 & 0.6698 \\
        \bottomrule
    \end{tabular}
\end{table*}

\begin{table*}[ht]
    \centering
    \tablefontsize
    \caption{Metallography-oriented supervised deep models (15) on \datasetname{} (macro over subsets).}
    \label{tab:deep-metal}
    \begin{tabular}{llccc}
        \toprule
        Model & Category & mIoU & Dice & Pixel Acc \\
        \midrule
        Metal-U-Net++ CLAHE (EfficientNet-B0) & Micrograph-contrast & \textbf{0.6427} & \textbf{0.7456} & \textbf{0.8684} \\
        Metal-MAnet RGB+Sobel (EfficientNet-B0) & Edge-aware & 0.5987 & 0.7059 & 0.8294 \\
        Metal-U-Net CLAHE (ResNet34) & Micrograph-contrast & 0.5857 & 0.6962 & 0.8144 \\
        Metal-U-Net RGB+Sobel (ResNet34) & Edge-aware & 0.5784 & 0.6860 & 0.7984 \\
        Metal-SegFormer CLAHE (MiT-B0) & Micrograph-contrast & 0.5766 & 0.6848 & 0.8258 \\
        Metal-U-Net Gray (ResNet34) & Micrograph-contrast & 0.5713 & 0.6905 & 0.7555 \\
        Metal-U-Net++ Gray (ResNet34) & Micrograph-contrast & 0.5646 & 0.6914 & 0.7667 \\
        Metal-DeepLabV3+ CLAHE (ResNet34) & Micrograph-contrast & 0.5600 & 0.6829 & 0.7603 \\
        Metal-LinkNet RGB+Sobel (ResNet34) & Edge-aware & 0.5362 & 0.6546 & 0.7630 \\
        MLography U-Net (2022-style) & Metallography-original & 0.5115 & 0.6347 & 0.7489 \\
        Metal-U-Net LBP Stack (ResNet34) & Texture-aware & 0.4882 & 0.5945 & 0.6861 \\
        Metal-SegFormer Gray (MiT-B2) & Micrograph-contrast & 0.4802 & 0.5961 & 0.6941 \\
        Metal-FPN Gabor Stack (ResNet34) & Texture-aware & 0.4508 & 0.5565 & 0.7075 \\
        Metal-U-Net Gabor Stack (ResNet34) & Texture-aware & 0.3900 & 0.5017 & 0.6471 \\
        Metal-UPerNet CLAHE (MiT-B2) & Micrograph-contrast & 0.3650 & 0.4664 & 0.6175 \\
        \bottomrule
    \end{tabular}
\end{table*}

\begin{table}[ht]
    \centering
    \tablefontsize
    \caption{Foundation/edge add-on models (6) on \datasetname{} (inference-only, macro over subsets).}
    \label{tab:foundation-edge}
    \begin{tabular}{lccc}
        \toprule
        Model & mIoU & Dice & Pixel Acc \\
        \midrule
        SAM ViT-Base (auto-mask) & \textbf{0.4565} & \textbf{0.5366} & \textbf{0.7546} \\
        SlimSAM-50 (auto-mask) & 0.4558 & 0.5361 & 0.7541 \\
        SlimSAM-77 (auto-mask) & 0.4490 & 0.5302 & 0.7515 \\
        TextureSAM-0.3 (auto-mask) & 0.4417 & 0.5265 & 0.7498 \\
        HED + Watershed & 0.3496 & 0.4082 & 0.7099 \\
        PidiNet + Watershed & 0.3474 & 0.3959 & 0.7210 \\
        \bottomrule
    \end{tabular}
\end{table}

\section{Benchmark Website Usage Guide (Appendix)}
The public benchmark companion website is available at \url{https://scientific-computing-user.github.io/amam-benchmark-site/}. It mirrors the same strict pair-only release used in this paper and provides an interactive view of subsets, masks, and benchmark result tables. Its role is to improve transparency and qualitative auditing (sample-level inspection, artifact-focused error review, and access to released result files), while the canonical evaluation protocol, splits, and reporting requirements remain those defined in the main paper. Table~\ref{tab:website_artifacts} maps each website artifact to its practical value for reproducible benchmarking submissions.

\begin{table}[ht]
    \centering
    \tablefontsize
    \caption{Companion website artifacts and their benchmark-submission relevance.}
    \label{tab:website_artifacts}
    \begin{tabular}{p{0.30\linewidth}p{0.64\linewidth}}
        \toprule
        Website artifact & Submission-relevant value \\
        \midrule
        Subset explorer (\texttt{index.html}) & Makes material/condition/magnification coverage explicit and auditable at sample level. \\
        Aligned image--mask browsing & Supports artifact-focused qualitative analysis (blur, pores, scratches, reflective confounds). \\
        Benchmark result report (\texttt{report.html}) & Exposes model-block comparisons with the same metric schema used in this manuscript. \\
        Released CSV bundles & Enables independent recomputation of tables/plots and direct baseline comparison for new methods. \\
        Repository zip + manifest & Provides a fixed release snapshot for reproducible submission packaging and protocol traceability. \\
        \bottomrule
    \end{tabular}
\end{table}

\section{Reproducibility Details (Appendix)}
This appendix documents how the quantitative tables were produced and how the representative prediction visualizations were generated. The classical 10-method survey used deterministic splits and seeds (seed 17), image resize to 256$\times$256, strict pair-only inclusion, and subset-wise Hungarian class alignment. The supervised deep extension used deterministic splits (seed 17), image size 192$\times$192, 5 epochs per model, and subset-macro reporting under a subset-unique global-class training head. The foundation/edge add-on block used deterministic inference and the same held-out split. Transformer variants used CPU fallback where required by MPS backend constraints. Full outputs are in \texttt{results/benchmark\_*.csv} (classical), \texttt{results/deep\_survey/deep\_*.csv} (supervised deep), and \texttt{results/foundation\_edge/foundation\_edge\_*.csv} (foundation/edge add-ons).

\section{Reporting and Validity Notes (Appendix)}
\textbf{Recommended reporting minimum.} For comparable submissions, we recommend publishing: (i) subset-macro mIoU as primary score, (ii) per-subset breakdown, (iii) train/validation/test manifest hash, (iv) at least three random seeds for supervised methods, and (v) preprocessing disclosure (e.g., CLAHE, edge channels, grayscale conversion). These low-overhead requirements make cross-paper comparisons materially more reliable.

\textbf{Known limitations.} This release is intentionally strict but still compact. Some subsets have small holdouts (e.g., high magnification cast iron), and label spaces remain subset-native rather than globally unified. We consider this a realistic benchmark constraint, not a defect: the purpose is to evaluate method stability under genuine material-regime variation rather than artificial label flattening.

\textbf{Threats to validity.} Even with deterministic splits and strict pair filtering, three risks remain for leaderboard claims. \emph{Release-version drift}: benchmark conclusions can shift if future releases silently alter tuple membership or label definitions. \emph{Subset-size imbalance}: very small held-out slices can increase ranking variance for methods with unstable class calibration. \emph{Semantic transfer ambiguity}: cross-subset transfer failures can mix two effects, representation limits and non-identical phase ontologies. For this reason, confidence intervals and per-subset confusion reporting should be treated as first-class outputs, not optional appendix material.

\textbf{Submission packet template.} To keep reported gains auditable, we recommend each benchmark submission include the following minimal artifacts:
\begin{enumerate}[leftmargin=*,nosep]
    \item immutable train/validation/test manifest hash and release tag;
    \item model architecture and initialization source (including pretrained checkpoint identity);
    \item preprocessing and augmentation configuration (e.g., CLAHE, edge channels, grayscale conversion);
    \item per-track metric table with confidence intervals over at least three random seeds;
    \item compute budget disclosure (GPU type, wall-clock training time, and total trainable parameters).
\end{enumerate}
These requirements are lightweight and primarily enforce comparability across submissions.

\end{document}
